\section{The local arithmetic source and a one-place criterion}

We establish the arithmetic and local inputs for the unrestricted
pro-$3$ theorem. They also give a separate criterion that assumes
local geometry only at $3$ and a ramification condition on the
\emph{whole section cocycle}. That criterion applies to every
Galois-stable quotient retaining geometric abelianization.

Put $K=\Q(\zeta_3)$ and $\lambda=1-\zeta_3$. For a finite set $S$ of
rational primes containing $3$, let $K_S(3)$ be the maximal pro-$3$
extension of $K$ unramified outside the places above $S$, and set
\[
 P_S=\Gal(K_S(3)/K),\qquad
 \Gamma_S=\Gal(K_S(3)/\Q).
\]
The second extension is Galois because its defining ramification
condition is invariant under $\Gal(K/\Q)$. We shall use the condition
\[
 H^1(P_S,\F_3)\longrightarrow H^1(K_\lambda,\F_3)
 \quad\hbox{is injective}.
 \tag{6.1}\label{eq:sourcecriterion}
\]

\begin{lemma}[Local source surjectivity]\label{lem:source}
If \eqref{eq:sourcecriterion} holds, the natural homomorphism
$G_{\Q_3}\to\Gamma_S$ is surjective. Condition
\eqref{eq:sourcecriterion} holds for
\[
 S=\{3\},\qquad S=\{3,7\},\qquad
 S=\{3,q\}\quad(q\text{ prime},\ q\equiv2\text{ or }5\pmod9).
\]
\end{lemma}
\begin{proof}
The image of $G_{K_\lambda}$ in $P_S$ is closed. If it were proper,
some finite $3$-group quotient would have proper image; a maximal
subgroup of that quotient has index three. The resulting nonzero
character $P_S\to\F_3$ would restrict to zero locally, contradicting
\eqref{eq:sourcecriterion}. Hence this local source surjects onto
$P_S$. The extension $K_\lambda/\Q_3$ has degree two, so the image of
$G_{\Q_3}$ also surjects onto $\Gamma_S/P_S$. It is all of $\Gamma_S$.

The imaginary quadratic field $K$ has class number one: its
Minkowski bound is $(2/\pi)\sqrt3<2$. Consequently Kummer theory
identifies $H^1(P_S,\F_3)$ with the group of $S$-units modulo cubes;
an element whose valuations outside $S$ are multiples of three can
have those valuations removed using principal ideals. Locally let
$U_i=1+\lambda^i\mathcal O_{K_\lambda}$. The groups of roots of unity,
the valuation, and the logarithm give
\[
 K_\lambda^\times/(K_\lambda^\times)^3
 =\langle\lambda\rangle\oplus\langle\zeta_3\rangle
   \oplus U_2/U_4.
\]
Here each of the first two factors has order three, and $U_2/U_4$
has dimension two over $\F_3$. Indeed $U_1=\mu_3\times U_2$, and
the logarithm identifies cubing on $U_2$ with multiplication by $3$,
whose $\lambda$-valuation is two. Thus $U_2^3=U_4$.

For $S=\{3\}$, the global basis is $\zeta_3,\lambda$, visibly
independent locally. If $q\equiv2$ or $5\pmod9$, then $q$ is inert
in $K$ and the global basis is $\zeta_3,\lambda,q$. Write
$c=(q+1)/3$, which is a $3$-adic unit. Since
$3=-\lambda^2/\zeta_3$,
\[
 -q=1-3c\equiv1+c\lambda^2+c\lambda^3\pmod{\lambda^4}.
\]
Its $U_2/U_4$ coordinate is nonzero, proving independence.
Finally, at $7$ put $\pi=2-\zeta_3=1+\lambda$; its norm is seven.
The global basis is $\zeta_3,\lambda,\pi,\bar\pi$.
Replacing the last element by $7=\pi\bar\pi$ preserves that basis.
The exact identity $\zeta_3\pi=1-\lambda^2$ and the congruence
\[
7=1-2\lambda^2/\zeta_3
 \equiv1+\lambda^2+\lambda^3\pmod{\lambda^4}
\]
give the two independent rows $(-1,0)$ and $(1,1)$ in $U_2/U_4$.
This proves \eqref{eq:sourcecriterion} in all the stated cases.
\end{proof}

\begin{lemma}[The compact cubic at three]\label{lem:cubiclocal}
For $E:v^2+v=u^3-7$, the group $E(\Q_3)$ is pro-$3$. Its
$3$-adic Kummer map is therefore injective. A global class in
$H^1(\Q,T_3E)$ that is Kummer at $3$ is the Kummer class of one of
the three points of $E(\Q)$.
\end{lemma}
\begin{proof}
The integral equation has discriminant $-3^9$, and is minimal since
its valuation is less than twelve. Its reduction has unique singular
point $(0,1)$ and three smooth $\F_3$-points, counting infinity.
Let $E_0$ be the inverse image of the smooth locus and $E_1$ the
formal subgroup. Thus $E_1$ is pro-$3$ and $E_0/E_1$ has order three.
The point $P=(3,4)$ has order three: its tangent has slope three,
and the addition formula gives $2P=(3,-5)=-P$.

We verify that every point with singular reduction belongs to
$P+E_0$ or $-P+E_0$. Such a point has integral coordinates
$u=3a,v=1+3b$. Its equation implies $b^2+b+1=3a^3$, so
$b=1+3w$. It follows that $a^3=1+3w+3w^2$. Cubes modulo nine
are $0,1,-1$, which forces $w\equiv0$ or $-1\pmod3$ and
$a\equiv1\pmod3$. Hence write
\[
 u=3+9z,\quad v=4+9w,\qquad
 w^2+w=3z+9z^2+9z^3.
\]
If $w\equiv0\pmod3$ and $z\ne0$, the slope in $Q-P$ is
$(w+1)/z$. If $z$ is a unit, the new $u$-coordinate is $1$ modulo
three and the reduction is smooth. If $z$ is not a unit, the new
coordinate has negative valuation and the sum reduces to infinity.
For $z=0$ this case gives $Q=P$. Thus $Q-P\in E_0$.
If $w\equiv-1\pmod3$, apply the same reasoning to $Q+P$,
whose slope is $w/z$; the case $z=0$ is $Q=-P$.
Consequently $E/E_0$ is generated by the nonzero class of $P$ and
has order three. The quotient $E/E_1$ has order nine, proving the
pro-$3$ assertion. The intersection of all $3^mE(\Q_3)$ is zero,
so the local Kummer map is injective.

The modular argument in the proof of Theorem~\ref{thm:solvable-nineteen-intro}
gives $E(\Q)\simeq\Z/3\Z$ and $T_3\Sha(\Q,E)=0$. At a finite
prime $v\ne3$, $E$ has good reduction and every class in
$H^1(\Q_v,T_3E)$ is Kummer. For example, the unramified cohomology
sequence has quotient $(T_3E(-1))^{\operatorname{Frob}_v}=0$ by
the Weil bounds, and its unramified term is the Kummer group of
the finite reduction. At the real place the cohomology is zero
because two is invertible on $T_3E$. Thus the hypothesized global
class lies in the inverse-limit $3$-Selmer group. The Kummer exact
sequence and $T_3\Sha=0$ identify that group with $E(\Q)$.
\end{proof}

\begin{lemma}[Action and change of tangent]\label{lem:action}
Let $n=3^r$, $r\geq1$, and put
$N=\pi_1((U_n)_{\bar\Q})^{(3)}$. Choose at a rational coordinate
cusp the tangent supplied by a coordinate parameter with leading
coefficient one. The resulting actual action on $N$ is unramified
outside $3$, and its restriction to $G_K$ has pro-$3$ image.
The relative cocycle between any two such standard rational
tangents is also unramified outside $3$.
\end{lemma}
\begin{proof}
The plane Fermat model and its coordinate boundary form a smooth
proper curve with an \emph{\'{e}tale} boundary over $\Z[1/3]$.
The chosen tangents extend as nonzero vectors there. At each
prime $v\ne3$, tame specialization identifies the pro-$3$
geometric covering category, including these tangential fiber
functors, with that of the geometric special fiber
\citep[Expos\'e XIII, \S2]{SGA1}. Inertia acts trivially on the
latter category and its fiber-functor isomorphisms. Therefore it
acts trivially on the actual based fundamental group and on every
path torsor between the stated tangents. This gives both
unramifiedness assertions, including the change-of-origin cocycle.
The integrality and nonvanishing of the tangent are essential in
this application of specialization.

There is also a direct pointed description using the triply punctured
line $B=\mathbb P^1\setminus\{0,-1,\infty\}$ with coordinate $s$.
The cover $U_n\to B$ is given by
\[
s=x^n,\qquad(-y)^n=1+s.
\]
If $F_0=\pi_1(B_{\bar\Q})^{(3)}$ with the unit tangent at zero,
then $F_0^{\ab}=\Z_3(1)^2$ and
\[
N_n=\ker(F_0\longrightarrow F_0^{\ab}/nF_0^{\ab}).
\]
This is an open characteristic subgroup. The equality with the full
pro-$3$ group of $U_n$ follows by taking the normal closure of a finite
$3$-cover under the $3$-group of deck transformations; its monodromy
is still a $3$-group. At $Q_0$ the explicit fiber functor has
$x=t$, $-y=(1+t^n)^{1/n}$ and $s=t^n$. Coefficientwise Galois action
on rational Puiseux series makes the arithmetic splittings compatible.
The other two standard cusps have the analogous coordinate description.
Thus the actual action on $N_n$ is the restriction of that on $F_0$,
and the inclusions $N_n\subset N_3\subset F_0$ are Galois stable.

On $F_0/\Phi(F_0)$ the action is scalar cyclotomic and is trivial
on $G_K$. The kernel of
$\operatorname{Aut}(F_0)\to\operatorname{Aut}(F_0/\Phi(F_0))$
is pro-$3$, by the finite $3$-group automorphism theorem and passage
to characteristic lower $3$-central quotients. Hence the actual
$G_K$-image on $F_0$, and on $N_n$, is pro-$3$.
For this punctured-line description the pointed unramification and
path-torsor statements above also follow directly from
\citet[Theorem 17.7 and Corollary 17.8, p.~51]{Wojtkowiak2009}.
\end{proof}

\Needspace{8\baselineskip}
\begin{theorem}[Restricted pro-$3$ Fermat sections]\label{thm:restricted3}
Let $n=3^r$, $r\geq1$, and fix one of the standard integral
tangents in Lemma~\ref{lem:action}. Let $M=N/J$, where $J$ is a
closed Galois-stable normal subgroup of $N$ contained in $[N,N]$;
the choice $J=1$ is permitted. Suppose $S$ is a finite set
containing $3$ and satisfies \eqref{eq:sourcecriterion}.

If a section of $M\rtimes G_\Q$ has its entire relative cocycle
trivial on inertia outside $S$ and its localization at $3$ is
point-or-cusp geometric, then it is globally cuspidal.
Within this ramification class its sections form exactly three
disjoint packets, one at each rational coordinate cusp, each
parametrized by
\[
 H^1(G_{\Q,S},\Z_3(1))
    =\bigoplus_{q\in S}\Z_3\,\kappa_3(q).
 \tag{6.2}\label{eq:restrictedpackets}
\]
In particular this holds for all the sets $S$ in
Lemma~\ref{lem:source}, proving Theorem~\ref{thm:restrictedintro}.
No local condition away from $3$ is needed
for the implication to cuspidality.
\end{theorem}
\begin{proof}
Push the section to the compact cubic using $X_n\to X_3\simeq E$.
This is possible through $M$ since that map has abelian geometric
target and $J\subset[N,N]$. Lemma~\ref{lem:cubiclocal} identifies
the resulting global compact class with a point of $E(\Q)$.
These three points are its coordinate cusps. The local center at
$3$ maps to precisely that point, since the compact $3$-adic
Kummer map is injective. Its inverse image in $X_n$ is supported
on the coordinate boundary. Moreover, $\Q_3$ has no nontrivial
odd-order roots of unity: its roots of unity are $\{1,-1\}$.
Thus this center is one of the three rational cusps of $X_n$,
say $Q$. The local section must be cuspidal at $Q$.

Let $\alpha$ be the actual action and $z$ the section cocycle.
Over $G_K$ the graph homomorphism
$g\mapsto(z(g),\alpha(g))$ has pro-$3$ image by
Lemma~\ref{lem:action}, and is unramified outside $S$. It factors
through $P_S$ by that group's defining maximality. Consequently
both $\alpha$ and $z$ factor through $\Gamma_S$. Replace the
origin by the standard tangent at $Q$. Its change-of-origin
cocycle is unramified outside $3$ by Lemma~\ref{lem:action}, so
the new action and cocycle still factor through $\Gamma_S$.

Choose a single element of $M$ conjugating the local section at
$3$ into the cusp group at $Q$, and conjugate the global section
by that same element. Its new cocycle is $I_Q$-valued on
$G_{\Q_3}$, where $I_Q\simeq\Z_3(1)$ denotes the peripheral
subgroup in $M$. Lemma~\ref{lem:source} says this local source
surjects onto $\Gamma_S$. Every value of the global cocycle is
therefore already a value occurring on the local source, and
lies in $I_Q$. This is exactly a global cusp section at $Q$.

For completeness the peripheral subgroup embeds in $M$ because
it embeds already in $N^{\ab}$ by the single boundary relation.
Its parameters are unique: filling $Q$ gives the pro-$3$ analogue
of \eqref{eq:boundary}, and the Tate module of the resulting
generalized Jacobian has zero invariants globally and over
$\Q_3$. Thus the peripheral maps on $H^1$ are injective. Distinct
centers are separated by the compact cubic Kummer classes at $3$.
Since the standard tangents are unramified outside $3$, a global
cusp parameter is unramified outside $S$ exactly when its
valuations away from $S$ vanish. Kummer theory over
$\Z[S^{-1}]$, whose ideal class group is zero, gives
\eqref{eq:restrictedpackets}. Every such cusp section satisfies
all local point-or-cusp conditions. This proves the classification.
\end{proof}

The whole-cocycle ramification hypothesis is essential to this
one-place argument. The unrestricted Selmer theorem below instead
applies source surjectivity to successive additive central differences;
it permits cusp parameters with infinite support.
